\documentclass[10pt,a4paper]{amsart}
\usepackage[margin=19mm]{geometry}
\usepackage{amsmath,amssymb,mathtools}
\usepackage[scr=rsfs,cal=euler]{mathalfa}
\usepackage{xfrac}
\usepackage[unicode,hidelinks,bookmarks=false]{hyperref}
\newcommand{\ie}{i.e.\ }
\newcommand{\cf}{cf.\ }
\newcommand{\resp}{respectively\ }
\newcommand{\Subsection}{\subsection}
\providecommand{\nobreakdash}{\nobreak\hbox{-}}
\setlength{\emergencystretch}{2em}
\multlinegap0pt
\newtheorem{comment}{Comment}
\usepackage{bm}
\usepackage{amssymb}
\usepackage{mathtools}
\usepackage[normalem]{ulem}
\usepackage{csquotes}
\usepackage[shortlabels]{enumitem}
\usepackage{xcolor}
\usepackage{tikz}
\newtheorem{proposition}{Proposition}
\newcommand{\Ti}{\textnormal{Ti}}
\newcommand{\scal}{\cdot} %minkowski product
\newcommand{\y}{\bm{y}} 
\title{Ti and Spi, Carrollian extended boundaries at timelike and spatial infinity}
\author{Jack Borthwick$^{\dagger}$, Maël Chantreau$^{\ddagger, \star}$, Yannick Herfray$^{\ddagger}$.}
\date{Source: arXiv:2412.15996v2 (CC BY 4.0)}
\begin{document}
\maketitle

\noindent\emph{Source and licence.} Ti and Spi, Carrollian extended boundaries at timelike and spatial infinity. Version arXiv:2412.15996v2 (CC BY 4.0), distributed by arXiv under the Creative Commons Attribution 4.0 International licence, http://creativecommons.org/licenses/by/4.0/, as recorded in the arXiv licence metadata for this exact version and shown on the abstract page https://arxiv.org/abs/2412.15996v2. Full licence notice: \texttt{licenses/arxiv-2412.15996-CC-BY-4.0.txt}.\par\smallskip

\noindent\textit{Editorial scope.} Counted unit: the article's proposition Proposition::MassiveScalar (source0.tex lines 429--442). The article's complete original proof of it is delivered with it (source0.tex lines 443--460, one \texttt{proof} environment of 18 lines). No mathematics is written, repaired or re-derived: the statement and the proof are the article's own lines, in the article's own notation, and no retained line is substituted or re-flowed.  Retained uncounted as the source-local context the proof needs: lines 407--409.  Nothing was omitted from any retained span, so the package carries no omission record.  No retained line contains a citation, so the wrapper carries no bibliography and no bibliography entry.  No retained line contains a figure environment, an image-inclusion command or drawing code, so no image is delivered and the package declares no asset dependency.  Editorial formatting only, in addition: an AMS \texttt{amsart} preamble that restates the article's own declarations and macros used by the retained lines, namely the environments \texttt{proposition} on separate counters, together with the standard packages those lines need.  The retained environments print in this document as Proposition 1 in order of appearance, whereas the article prints the same items with the numbering of its own full document.  The complete native source of the article as distributed in the arXiv e-print for this exact version is bundled unchanged as \texttt{sources/170485/source0.tex}.
\begin{equation}\label{Massive scalar field: Fourier coefficients}
    \Phi(X) = \frac{m^{n-1}}{2 (2\pi)^n}\int \frac{d^{n}\bm{k}}{\sqrt{1 + \bm{k}^2}} \Big( a(\bm{k}) e^{i m \tilde{P} \scal X} + a(\bm{k})^{\dagger} e^{-i m \tilde{P} \scal X}\Big).
\end{equation}

\begin{proposition}\mbox{}

\label{Proposition::MassiveScalar}

Let $(u, \y )$ be a point of $\Ti^{+}$, then the limit
\begin{align}
    \varphi(u, \y) &:= \lim_{\rho\to 0} \quad \frac{1}{2}\Big(    \Phi(X) - \frac{1}{i m} \nabla^a \rho_u \nabla_a\Phi(X) \Big) \rho_u^{-\frac{n}{2}}e^{in\frac{\pi}{4}} e^{i m \rho_u^{-1}} \label{Eq::AsymptoticScalar1}
\end{align}
exists and defines a field on $\Ti^+$. In terms of the Fourier coefficients \eqref{Massive scalar field: Fourier coefficients}, this reads
\begin{equation}
    \varphi(u, \y) = \, \frac{m^{\frac{n}{2}-1}}{2 (2\pi)^{\frac{n}{2}}}  a(\y) e^{-im u}. \label{Eq::AsymptoticScalar2}
\end{equation}
We will call this function the scattering data induced by the massive field on $\Ti$.
\end{proposition}

\begin{proof}\mbox{}
\begin{comment}
First, the symbol $(0^+-i\rho_u)^{-\frac{n}{2}}$ is shorthand for: $\displaystyle \lim_{\varepsilon \to 0^+} e^{-\frac{n}{2}\log(\varepsilon-i\rho_u)}$ and $\log$ is the principal branch of the complex logarithm, i.e. \[ (0^+-i\rho_u)^{-\frac{n}{2}}=\rho_u^{-\frac{n}{2}}e^{in\frac{\pi}{4}}.\]
\end{comment}
By hypothesis, $\nabla^a \rho_u = \rho^{2}\left( - \partial_{\rho} + h^{\alpha\beta}\nabla_{\alpha} u \,\partial_{\beta} \right) +O(\rho^3)$; note that the index has been raised using the rescaled inverse metric $g^{ab}$. Consequently,
\begin{equation}
    - \frac{1}{i m}\nabla^a \rho_u \nabla_a\Phi(X) \underset{\rho \to 0}{\sim} \rho^{\frac{n}{2}} \, \frac{m^{\frac{n}{2}-1}}{2 (2\pi)^{\frac{n}{2}}}  \left( a(\y) e^{- im \rho^{-1}-in\frac{\pi}{4}} - \, a(\y)^{\dagger} e^{im \rho^{-1}+in\frac{\pi}{4}} \right)  + O(\rho^{\frac{n}{2}+1} )
\end{equation}
from which we obtain
\begin{equation}
\frac{1}{2}\left(\Phi(X) - \frac{1}{i m}\nabla^a \rho_u \nabla_a\Phi(X) \right) \underset{\rho \to 0}{\sim} \rho^{\frac{n}{2}} \, \frac{m^{\frac{n}{2}-1}}{2 (2\pi)^{\frac{n}{2}}}  \, a(\y) e^{- im \rho^{-1}-in\frac{\pi}{4}}    + O(\rho^{\frac{n}{2}+1} ).
\end{equation}
This selects one of the asymptotic plane waves, and the remaining factors lift it to a field on $\Ti^+$:
\begin{equation}
\frac{1}{2}\left(\Phi(X) - \frac{1}{i m}\nabla^a \rho_u \nabla_a\Phi(X) \right)  (0^+-i\rho_u^{-\frac{n}{2}})  e^{i m  \rho_u^{-1}} \underset{\rho \to 0}{\sim}  \frac{m^{\frac{n}{2}-1}}{2 (2\pi)^{\frac{n}{2}}} \, a(\y) e^{- im u}    + O(\rho).
\end{equation}

\end{proof}

\end{document}
